\chapter{Introduction}
\label{ch:introduction}

% ================================================================
% 1.1 Context and motivation
% ================================================================
\section{Context and motivation}
\label{sec:context-motivation}

\subsection{Growth and evolution of network threats}
\label{sec:growth-evolution-threats}

The volume and variety of attacks on network infrastructure keep growing. Two
independent sources of evidence, gathered from very different vantage
points, point the same way.

The first is regulatory and incident-driven. The 2026 edition of the
\gls{enisa} \emph{Threat Landscape} report analyzes 8{,}257 incidents
recorded between January and December 2025. They come mainly from open
sources, plus anonymised information shared by Member States and partners.
\gls{ddos} was the primary incident type at 51.3\%, followed by
unauthorised access at 39.5\% \citep{enisa2026threatlandscape}. The report
calls these \gls{ddos} attacks low-impact. It also warns that ransomware and
\gls{ddos} are more immediately visible than other threats, because their
operators often claim them directly, so the counts should be read with this
in mind. By assessed objective, 57.3\% of the incidents were ideology-driven,
particularly hacktivist \gls{ddos} campaigns, and about 29\% were financially
motivated. The report states that financially motivated activity remains the
most impactful threat in the short term.
Among social engineering techniques, phishing made up 77.8\%. The report
could identify an intrusion vector for only 5.2\% of the unauthorised access
incidents, and in 60.4\% of those the vector was a vulnerability. The report
also observes a growing use of AI by attackers, mostly consumer-grade tools
that augment existing skills and do not give breakthrough capabilities.

The second source comes from a commercial mitigation provider, which
observes attack traffic directly. Its figures are vendor data and cannot be
checked independently. For 2025, StormWall reports that the total number of
\gls{ddos} attacks grew by 198\% and that about 70\% of websites
experienced at least one attack during the year
\citep{stormwall2026forecast}. The character of the attacks is changing as
well. The same source reports that multi-vector attacks, which combine two
or more vectors in one incident, grew by 83\% and appeared in nearly one
third of all incidents. It forecasts a share of up to 65\% for 2026. It also
describes a 4.8~Tbps attack of the carpet-bombing type, which spreads
malicious traffic across IP ranges and subnets instead of concentrating it
on a single vector. Reconnaissance attacks are a further case. They usually
last up to 15 minutes at relatively low intensity, often stay below
filtering thresholds, and made up 25\% of the recorded incidents. This
matters for the thesis because attacks of this kind are hard for classical
volumetric thresholds to catch, a point taken up again with the threshold
defences of \cref{sec:threshold-rate-limiting}.

Two conclusions follow. First, the attacker population is heterogeneous and
changes over time. Hacktivist volumetric noise, financially motivated
application-layer abuse and slow, targeted reconnaissance coexist on the
same networks, so a detector tuned for one profile will misjudge the
others. Second, the growth itself keeps moving. With increases as large as
the 198\% reported above, any approach whose maintenance cost grows with the
number of known attack signatures is bound to fall behind. These two points
lead to the two limitations examined next. Purely signature-based detection
cannot keep pace, and purely data-driven detection, although it addresses
that problem, brings failure modes of its own.
% Note: the ENISA figures are from the 2026 edition (September 2026). The
% report gives the financially motivated share as 29.2% in Chapter 1 and as
% 29.3% in the Executive Summary, hence "about 29%" in the text.

\subsection{Limits of signature-based detection}
\label{sec:signature-not-enough}

The paradigm that has dominated operational \gls{nids} deployments since the
1990s is \emph{misuse}, or \emph{signature-based}, detection. NIST defines a
signature as ``a pattern that corresponds to a known threat'', and
signature-based detection as ``the process of comparing signatures against
observed events to identify possible incidents''
\citep{scarfone2007nist80094}. Snort, whose original design was published by
\citet{roesch1999snort}, is the best-known tool of this kind. The approach
is also simple. The NIST guide calls it ``the simplest detection method'',
because it only compares the current unit of activity, such as a packet or a
log entry, to a list of signatures using string comparison. The same guide
describes it as ``very effective at detecting known threats''
\citep{scarfone2007nist80094}.

The guide is explicit about the opposite weakness. Signature-based detection
is ``largely ineffective at detecting previously unknown threats, threats
disguised by the use of evasion techniques, and many variants of known
threats''. Such tools also have ``little understanding of many network or
application protocols and cannot track and understand the state of complex
communications'' \citep{scarfone2007nist80094}. \citet{khraisat2019survey}
make the same point. Signature-based systems give excellent detection
accuracy on previously known intrusions. They have difficulty with zero-day
attacks, because no matching signature exists in the database until the
signature of the new attack has been extracted and stored. The survey also
lists the need to update the signature set frequently, and the poor fit for
multi-step attacks, among the disadvantages.

Here the growth documented above turns a known limitation into an
operational problem. A signature database has to be maintained. Every new
attack variant that should be caught needs an analyst, or an automated
pipeline standing in for one, to write, test and deploy a new rule, and
every added rule raises the matching cost paid by all later traffic. If
attack volumes grow as fast as the industry figures suggest, and attackers
keep diversifying their tactics, the rule base must grow with them. The part
of that growth that is genuinely new, and not a cosmetic variant of
something already covered, is the part a signature-based system cannot
address until a human closes the gap. The result is a delay between the
appearance of an attack technique and the system's ability to recognize it,
the \emph{zero-day gap}. Its length depends on analyst response time and not
on anything the detector can improve by itself. This gap motivated the move
to learning-based detection, which trades it for a different set of
problems.

\subsection{Limits of data-driven detection}
\label{sec:datadriven-not-enough}

The limits of misuse detection led, from the late 1990s and with growing
force after the arrival of deep learning, to \gls{ml}- and \gls{dl}-based
detection. A common framing is \emph{anomaly detection}. Instead of matching
traffic against known-bad patterns, a model is trained on normal behavior and
flags deviations from it. The promise is to catch attacks never seen before, which would close the
zero-day gap. In practice, the literature has spent over a decade showing
why this is harder than the successes of machine learning in image
recognition, speech or language suggest.

The foundational critique is that of \citet{sommer2010outside}. They observe
that operational deployments use almost only misuse detectors. They argue
that finding attacks is fundamentally different from the tasks where machine
learning works well, such as product recommendation, optical character
recognition or spam detection. A classifier needs training examples of every
class. An anomaly detector aimed at novel attacks cannot train on the
attacks of interest and sees only normal traffic. The authors call this the
\emph{closed world assumption}, following Witten et al. They add further
difficulties: a very high cost of errors, a semantic gap between a detected
deviation and a report an operator can act on, a diversity of network
traffic that makes it hard to find a stable notion of normality, and an
evaluation that is harder than building the detector. They argue mostly at
an intuitive level and not formally. This matters for the growth pattern of
the previous subsections. If traffic and attacks keep shifting, a model
fitted to yesterday's data gives no guarantee about tomorrow's.

A second line of critique concerns how the field measures whether a
detector works at all. \citet{arp2022dosdonts} identify ten pitfalls of
machine learning in security. They review thirty papers from four top
security conferences, on topics from malware detection to network intrusion
detection. Every paper suffers from at least three of the pitfalls. Four
of them matter most here. \emph{Data snooping} lets information that is not
available in practice reach training, including test data and data from the
future. \emph{Spurious correlations} lead a model to learn artifacts
unrelated to the security problem instead of the attack behavior. The
\emph{base rate fallacy} ignores a large class imbalance when interpreting
performance. \emph{Lab-only evaluation} tests a system in a closed-world
setting with limited diversity, without considering that data change over
time. In four case studies the authors show that individual pitfalls can
lead to unrealistic performance. The base-rate problem is especially serious
for network traffic, where malicious flows are a small fraction of the
total. The authors give an example with a class ratio of 1:100. A detector
with 99\% true positives and 1\% false positives then produces 100 false
positives for every 99 true positives. The ratio in real traffic can be more
extreme, and the effect grows with it. This is worked out numerically in
\cref{sec:evaluation-practice}, where the benchmarking protocol of this
thesis is motivated.

Taken together, the two critiques show that purely data-driven detection
does not just swap one problem for the opposite one. It exchanges a
well-understood if inconvenient limitation, blindness to the unknown, for
harder ones: an unstable notion of normality, and an evaluation practice
that can make brittle models look robust. More model capacity or more
training data fixes neither, because both problems concern what is learned
and how success is measured, not the expressive power of the learner. This
gap motivates injecting explicit domain knowledge into the learning process,
which is the neurosymbolic approach studied in this thesis. Detection quality
is then not treated as something that more data or a bigger network will
eventually solve. \Cref{ch:learning-paradigms,ch:ltn} develop the argument
once the machine-learning and logic background is in place.

% ================================================================
% 1.2 Problem statement
% ================================================================
\section{Problem statement}
\label{sec:problem-statement}

The argument so far leads to an uncomfortable observation. Neurosymbolic
integration is the kind of remedy the previous section calls for. It injects
the explicit domain knowledge a purely data-driven detector lacks, without
giving up the pattern-recognition capacity that made data-driven detection
attractive. The literature accordingly proposes it as a way to improve
generalization, handle rare or unseen attack classes and produce explainable
alerts. Yet this thesis does not set out to show that neurosymbolic
detection works. It asks \emph{whether}, and under which conditions, it is
worth adopting at all, because the available evidence does not support a
confident answer either way.

The difficulty is not a lack of individual results but the lack of a shared
methodology to produce and compare them. The studies surveyed in
\cref{ch:related-work} evaluate neurosymbolic detectors under different
protocols: different datasets, different splits, different sets of encoded
rules, and mostly one or two datasets per paper. It stays open whether a
reported gain comes from the approach or from the benchmark it was measured
on. Reported outcomes are mostly limited to accuracy or closely related
classification metrics, with computational cost, such as training time, inference latency or memory
footprint, reported by few papers, and almost never under
the throughput and latency limits of a real deployment. Because the
studies rarely share a protocol, their results can hardly be aggregated or
contrasted. The
field collects isolated positive findings and has no comparable picture of
when the approach helps and when it does not.

This gap has a practical cost. A security operations engineer who must
decide whether to adopt a neurosymbolic detector, instead of or next to an
existing machine-learning pipeline, is asking a concrete question: under
which conditions is a neurosymbolic detector worth its extra complexity? The
literature offers no shared answer. It offers data points, each obtained
under its own assumptions, but no systematic account of the operating regime.
That regime is the combination of data availability, class imbalance,
exposure to unseen attacks and knowledge quality in which the added
engineering complexity gives a measurable advantage over a matched, purely
neural baseline, and the regime in which it does not. This thesis supplies
that account through a benchmark designed for comparability and not for
showcasing a single best result. \Cref{sec:research-questions,sec:contributions}
turn it into research questions and contributions.

% ================================================================
% 1.3 Research questions
% ================================================================
\section{Research questions}
\label{sec:research-questions}

The problem statement is split into five research questions. The first four
cover the three axes a fair comparison needs, detection quality, cost and
usability, plus the operating-regime question that ties them together. The
fifth makes the exercise itself a methodological deliverable, because
RQ1--RQ4 can be answered credibly only once a benchmark protocol able to
support the answer exists.

\begin{description}
  \item[RQ1 -- Detection quality.] Does injecting domain knowledge through
    a \gls{ltn} improve detection over an identical, purely neural model,
    and on which classes or attack families does the improvement, or its
    absence, concentrate? This is the most direct test of the neurosymbolic
    claim. With architecture, features and training budget held fixed, a
    difference in detection quality can be attributed to the logic-based
    knowledge component and not to model capacity. Reporting the answer per
    attack family avoids merging different failure modes, such as volumetric
    \gls{ddos} and rare user-to-root attacks, into one aggregate number.

  \item[RQ2 -- Operating regime.] Under which conditions does the
    neurosymbolic advantage of RQ1, if any, appear and disappear? The
    conditions considered are the amount of labelled data, the degree of
    class imbalance, exposure to attacks unseen in training, shift across
    datasets and the quality of the encoded knowledge. This is the question
    the problem statement finds missing in the literature. Individual papers
    report a result under one set of conditions, and no systematic account
    exists of the boundary between the regime where the approach helps and
    the regime where it does not, including the possibility that it never
    helps.

  \item[RQ3 -- Cost.] What is the computational price of the neurosymbolic
    approach in training and in inference, in batch and in streaming
    operation? A detector that improves quality at an unreported or
    prohibitive cost is of no use to an operator with a latency or
    throughput budget. Cost is therefore an independently reported axis and
    not a footnote to a quality result.

  \item[RQ4 -- Usability.] What does adopting the neurosymbolic approach
    cost in engineering effort, to build, validate and maintain the
    knowledge base, and what does it return in alerts that a human analyst
    can interpret and audit? Quality and computational cost do not show
    whether an organization can realistically build, maintain and trust the
    system. This question states that dimension explicitly and does not
    leave it to informal claims about interpretability.

  \item[RQ5 -- Methodology.] What must a fair, reproducible benchmark for
    machine-learning versus neurosymbolic network traffic analysis look
    like? RQ1--RQ4 can be answered credibly only when matched baselines,
    controlled confounds and comparable cost and usability measurements are
    defined well enough for another researcher to reproduce them. RQ5 is
    both a question in its own right and a precondition for the other four.
    Its answer is the methodology of \cref{ch:benchmarking-methodology}.
\end{description}

% ================================================================
% 1.4 Contributions
% ================================================================
\section{Contributions}
\label{sec:contributions}

Answering the research questions yields five contributions. The first is the
methodological instrument that makes the others possible. The next three are
the empirical and engineering results obtained with it. The last introduces
evaluation signals that this thesis argues the mainstream \gls{nids}
literature is missing.

\begin{description}
  \item[C1 -- A reproducible benchmark suite.] A benchmark suite for
    comparing \gls{ml} and neurosymbolic network traffic analysis along
    three independent axes, detection quality, computational performance
    and operational usability, in both batch and real-time settings. It
    answers RQ5 and is the foundation of C2 to C4. It is developed in
    \cref{ch:benchmarking-methodology} and instantiated in
    \cref{ch:implementation}.

  \item[C2 -- A systematic empirical comparison.] A controlled,
    matched-baseline comparison of a \gls{ltn}-based detector against
    \gls{ml}/\gls{dl} baselines, repeated across several datasets and
    several stress scenarios: scarce labels, class imbalance, unseen
    attacks, cross-dataset transfer and deliberately degraded domain
    knowledge. It answers RQ1 and supplies the evidence from which RQ2 is
    answered. The results are reported in \cref{ch:results}.

  \item[C3 -- A characterisation of the operating regime.] An explicit
    account of the conditions under which neurosymbolic integration pays
    off, built from the stress-scenario results of C2, including the
    negative ones. It answers RQ2 and is reported with the same rigor
    whether the outcome favors the neurosymbolic approach or not, since a
    negative result is as useful to a practitioner as a positive one.

  \item[C4 -- A streaming reference pipeline.] A reference implementation
    that runs an \gls{ltn}-trained detector on live, replayed traffic and
    reports end-to-end detection latency and throughput, not only offline
    batch accuracy. It supplies the real-time half of the answer to RQ3.
    The batch half comes from the cost measurements of C1 on the models
    trained for C2. The pipeline is described in \cref{sec:realtime-pipeline}.

  \item[C5 -- Two NeSy-specific evaluation metrics.] Two metrics that the
    mainstream \gls{nids} literature does not report but that follow
    naturally once a logic-based knowledge component is part of the system.
    \emph{Axiom-level satisfaction} is an explanation signal: it shows how
    well the predictions satisfy each individual domain rule and not only
    the aggregate label. The \emph{domain-rule violation rate} measures how
    often the predictions themselves break a domain rule. The first
    contributes to the answer to RQ4, giving the usability axis of C1 a
    quantitative instrument, and the second adds to detection quality. Both
    are defined formally in \cref{sec:detection-quality-dimension} before
    being measured.
\end{description}

Contributions C1 to C5 are set against the research questions in
\cref{ch:conclusions}, once the evidence for each has been presented.

% ================================================================
% 1.5 Methodology in brief
% ================================================================
\section{Methodology in brief}
\label{sec:methodology-brief}

The full methodology is developed in \cref{ch:benchmarking-methodology} and
executed in \cref{ch:implementation,ch:results}. Here it is only outlined,
without anticipating the design choices and their justification.

The center of the thesis is a comparison between a \gls{ltn}-based detector
and a set of \gls{ml}/\gls{dl} baselines matched to it as closely as the two
paradigms allow. They share the input features and the preprocessing
pipeline, and the baseline uses a neural architecture identical to the one
embedded in the \gls{ltn} formulation. A measured difference in detection
quality can then be attributed to the logic-based domain knowledge and not
to a difference in model capacity or data access. The comparison runs on
established \gls{nids} datasets and is repeated under stress scenarios:
fewer labels, stronger class imbalance, attack families withheld from
training, transfer across datasets, and domain knowledge that is deliberately
degraded or partly wrong. Each scenario targets one of the conditions named
in RQ2.

Three evaluation dimensions are measured for every model in every scenario,
not accuracy alone. \emph{Detection quality} uses a metric set that stays
informative under severe class imbalance and is reported together with the
two neurosymbolic signals of C5. \emph{Computational cost} covers training
time and inference latency and throughput, in a batch setting and in a
streaming pipeline that replays traffic and measures end-to-end detection
delay. \emph{Usability} is a rubric-based account of the engineering effort
needed to build and maintain the knowledge base, set against the
interpretability the resulting alerts provide. Differences between models
are assessed with paired statistical tests on shared data splits. Every
experiment is run with several random seeds, and the configuration,
environment and code are released so that a third party can reproduce any
result reported in \cref{ch:results}.

% ================================================================
% 1.6 Scope and limitations
% ================================================================
\section{Scope and limitations}
\label{sec:scope-limitations}

The methodology is deliberately bounded, to keep the comparison controlled
and to fit the scope of a single-author thesis. Four exclusions apply
throughout. They are stated here so that the results are read with the right
boundaries in mind.

First, the thesis does not address inspection of encrypted payload or TLS
fingerprinting, and it does not inspect payload at all. On the flow datasets
the detectors use flow-level and header-level features of the kind
introduced in \cref{ch:foundations}. On NSL-KDD they also use the content
features that come with its records. These features were computed when the
dataset was built, and no flow record carries them. Second, adversarial
machine learning is considered only through the simple evasion perturbations
of one stress scenario, experiment E10 in \cref{ch:benchmarking-methodology}. A
systematic study of robustness against an adaptive attacker is left to
future work, discussed in \cref{sec:future-work}. Third, the runs that train a
model run on a server, with one exception, and the evaluations of saved models and
the timings of E13 run on a laptop (\cref{sec:impl-environment}). Distributed or
multi-node training, and the scalability questions that come with it, are out of
scope.
Fourth, the thesis stops at the detection and measurement layer. Production
integration with a security operations center, meaning SIEM ingestion, alert
ticketing and automated response, is neither implemented nor evaluated,
although the usability findings of RQ4 are meant to inform that context.

These boundaries limit what \cref{ch:results,ch:conclusions} can claim. The
comparison targets what a neurosymbolic detector costs and delivers as a
detection component, evaluated on realistic but decrypted, single-node,
non-adversarial traffic. It makes no claim about behavior under encrypted
payloads, an adaptive adversary, distributed deployment or full SOC
integration, each of which would need its own study.

% ================================================================
% 1.7 Structure of the thesis
% ================================================================
\section{Structure of the thesis}
\label{sec:thesis-structure}

The remaining ten chapters form three parts. The first, background and
design, comprises
\cref{ch:foundations,ch:attacks-ddos,ch:learning-paradigms,ch:ltn,ch:related-work,ch:usecases-design}.
The second, benchmarking and implementation, comprises
\cref{ch:benchmarking-methodology,ch:implementation}. The third, results and
conclusions, comprises \cref{ch:results,ch:conclusions}.

\Cref{ch:foundations} sets the vocabulary used in the rest of the thesis:
what is observed in network traffic and at which layer, how it is acquired
and turned into features, the classical detection approaches, the relevant
standards and taxonomies, and the recurring flaws in evaluation, treated
there in more quantitative depth.
\Cref{ch:attacks-ddos} develops the attack taxonomy, with emphasis on
\gls{ddos} given its weight in the incident figures above. It closes by
translating informal attacker knowledge into candidate logical axioms, which
\cref{ch:ltn} later formalizes.
\Cref{ch:learning-paradigms} surveys machine learning, deep learning and
symbolic AI on their own terms, then introduces the taxonomy of
neurosymbolic integration strategies in which the \gls{ltn} approach of this
thesis is placed.
\Cref{ch:ltn} presents \gls{ltn} in full, with its logical syntax, tensor
grounding, fuzzy semantics and training as satisfiability maximization, and
formalizes the network-security knowledge base used in the experiments.
\Cref{ch:related-work} surveys prior work on \gls{ml}/\gls{dl} for intrusion
detection and on neurosymbolic AI in cybersecurity, and places the research
gap of \cref{sec:problem-statement} against it.
\Cref{ch:usecases-design} turns the research questions into use cases,
functional and non-functional requirements, and the reference system
architecture, batch and streaming, used in the rest of the thesis.
\Cref{ch:benchmarking-methodology} is the methodological core announced in
\cref{sec:methodology-brief}. It gives the full benchmark design answering
RQ5 and delivering C1, including the three evaluation dimensions, the
statistical analysis plan and the threats to validity.
\Cref{ch:implementation} describes how the methodology is implemented: data
pipelines, models under test, the real-time pipeline delivering C4, and the
experimental setup.
\Cref{ch:results} reports and discusses the results against RQ1 to RQ4,
including the stress-scenario findings behind C2 and C3 and the usability
and explainability measurements behind C5.
\Cref{ch:conclusions} answers RQ1 to RQ5 explicitly, revisits C1 to C5,
distills a practical decision checklist for practitioners and outlines the
future work mentioned in \cref{sec:scope-limitations}.